\clearpage
\SourceLeaf{018}\pdfbookmark[0]{Inhalt.}{contenta}
\begin{center}
{\Large INHALT.}\par\bigskip
{\large ERSTES BUCH.}\par\medskip
Allgemeine Beziehungen zwischen den Größen, durch die die Bewegungen der Himmelskörper um die Sonne bestimmt werden.
\end{center}
\noindent\begin{tabular}{@{}p{.18\linewidth}p{.78\linewidth}@{}}
Abschnitt I. & \emph{Beziehungen für einen einzelnen Ort auf der Bahn.}\\[.4em]
Abschnitt II. & \emph{Beziehungen für einen einzelnen Ort im Raum.}\\[.4em]
Abschnitt III. & \emph{Beziehungen zwischen mehreren Orten auf der Bahn.}\\[.4em]
Abschnitt IV. & \emph{Beziehungen zwischen mehreren Orten im Raum.}
\end{tabular}
\begin{center}\bigskip
{\large ZWEITES BUCH.}\par\medskip
Untersuchung der Bahnen der Himmelskörper anhand geozentrischer Beobachtungen.
\end{center}
\noindent\begin{tabular}{@{}p{.18\linewidth}p{.78\linewidth}@{}}
Abschnitt I. & \emph{Bestimmung der Bahn aus drei vollständigen Beobachtungen.}\\[.4em]
Abschnitt II. & \emph{Bestimmung der Bahn aus vier Beobachtungen, von denen nur zwei vollständig sind.}\\[.4em]
Abschnitt III. & \emph{Bestimmung der Bahn, die mit beliebig vielen Beobachtungen möglichst gut übereinstimmt.}\\[.4em]
Abschnitt IV. & \emph{Über die Bahnbestimmung unter Berücksichtigung von Störungen.}
\end{tabular}
\par\bigskip\noindent\textsc{Tafeln.}
